\documentclass[10pt]{article}

% ---------- Document Identity (update these only) ----------
\newcommand{\DocStage}{}
\newcommand{\DocTitle}{Your Road to Tier One SOF}
\newcommand{\ProgramName}{182-Week Civilian-to-Selection Readiness Pipeline}
\newcommand{\ProgramSubtitle}{}

% ---------- Page ----------
\usepackage[legalpaper,left=0.5in,right=0.5in,top=0.6in,bottom=0.45in]{geometry}

% ---------- Typography ----------
\usepackage[T1]{fontenc}
\usepackage[utf8]{inputenc}
\usepackage{libertinus}
\usepackage{microtype}
\usepackage{parskip}
\usepackage{ragged2e}
\usepackage{textcomp}
\AtBeginDocument{\color{black}}
\AtBeginDocument{\pagecolor{white}}
\AtBeginDocument{\justifying}

% ---------- Landscape pages ----------
\usepackage{pdflscape}

% ---------- Color ----------
\usepackage[table]{xcolor}
\definecolor{StageBlue}{HTML}{1F4E79}
\definecolor{StageBlueDark}{HTML}{163A5A}
\definecolor{LightGray}{HTML}{F2F4F7}
\definecolor{MidGray}{HTML}{D9DEE5}
\definecolor{RuleGray}{HTML}{C7CED8}
\definecolor{HeaderInk}{HTML}{000000}
\definecolor{Ink}{HTML}{000000}

% ---------- Utility ----------
\usepackage{enumitem}
\sloppy
\setlength{\emergencystretch}{2em}

% ---------- Boxes / UI ----------
\usepackage[most]{tcolorbox}

\tcbset{
  charterTitle/.style={
    enhanced,
    colback=StageBlue!4,
    colframe=StageBlue,
    boxrule=1.25pt,
    arc=3pt,
    left=16pt,right=16pt,top=14pt,bottom=14pt,
    borderline north={3.2pt}{0pt}{StageBlue},
    borderline west={4pt}{0pt}{StageBlue},
    borderline south={1.25pt}{0pt}{StageBlue},
    drop shadow={black!25!white},
  },
  charterRule/.style={
    enhanced,
    colback=LightGray,
    colframe=RuleGray,
    boxrule=0.85pt,
    arc=3pt,
    left=12pt,right=12pt,top=10pt,bottom=10pt,
    borderline west={3pt}{0pt}{StageBlue},
  },
  charterBadge/.style={
    enhanced,
    colback=StageBlue,
    coltext=white,
    colframe=StageBlueDark,
    boxrule=0.6pt,
    arc=2pt,
    left=6pt,right=6pt,top=3pt,bottom=3pt,
  }
}
\newtcbox{\Badge}{nobeforeafter,charterBadge}

% ---------- Lists ----------
\setlist[itemize]{leftmargin=1.5em, itemsep=0.25em, topsep=0.25em}
\setlist[enumerate]{leftmargin=1.8em, itemsep=0.25em, topsep=0.25em}

% ---------- TOC ----------
\usepackage{tocloft}
\setcounter{tocdepth}{3}
\setlength{\cftbeforesecskip}{3pt}
\setlength{\cftbeforesubsecskip}{2pt}
\setlength{\cftbeforesubsubsecskip}{0pt}
\renewcommand{\cftsecfont}{\bfseries}
\renewcommand{\cftsecpagefont}{\bfseries}
\renewcommand{\cftsubsecfont}{\normalfont}
\renewcommand{\cftsubsecpagefont}{\normalfont}
\renewcommand{\cftsubsubsecfont}{\normalfont}
\renewcommand{\cftsubsubsecpagefont}{\normalfont}
\setlength{\cftsecindent}{0pt}
\setlength{\cftsubsecindent}{1.25em}
\setlength{\cftsubsubsecindent}{2.8em}
\setlength{\cftsecnumwidth}{2.1em}
\setlength{\cftsubsecnumwidth}{2.8em}
\setlength{\cftsubsubsecnumwidth}{3.6em}

% Remove the built-in TOC heading so only the custom heading is shown
\makeatletter
\renewcommand{\tableofcontents}{\@starttoc{toc}}
\makeatother

% ---------- Header/Footer: page number only ----------
\usepackage{fancyhdr}
\pagestyle{fancy}
\fancyhf{}
\renewcommand{\headrulewidth}{0pt}
\renewcommand{\footrulewidth}{0pt}
\fancyfoot[C]{\thepage}

% ---------- Section formatting ----------
\usepackage{titlesec}

\titleformat{\section}
  {\Large\bfseries\color{Ink}}
  {\thesection}{0.6em}{}
  [\vspace{0.25em}\textcolor{StageBlue}{\rule{\linewidth}{0.8pt}}]

\titleformat{\subsection}
  {\large\bfseries\color{Ink}}
  {\thesubsection}{0.6em}{}

\titleformat{\subsubsection}
  {\normalsize\bfseries\color{Ink}}
  {\thesubsubsection}{0.6em}{}

\titlespacing*{\section}{0pt}{0.95em}{0.35em}
\titlespacing*{\subsection}{0pt}{0.65em}{0.25em}
\titlespacing*{\subsubsection}{0pt}{0.55em}{0.2em}

% ---------- Table engine ----------
\usepackage{tabularray}
\UseTblrLibrary{booktabs}

% ---------- tabularray: GLOBAL table treatment ----------
\SetTblrInner{
  cells = {fg=black, bg=white, valign=t},
}

% ---------- Header helpers ----------
\newcommand{\Hdr}[1]{\shortstack[c]{\strut #1\strut}}

% ---------- Lines ----------
\newcommand{\TitleLine}{\textcolor{StageBlue}{\rule{0.98\linewidth}{0.95pt}}}
\newcommand{\SubTitleLine}{\textcolor{RuleGray}{\rule{0.98\linewidth}{0.65pt}}}

% ---------- Continued on next page ----------
\newcommand{\ContinuedOnNextPageInline}{%
  \par
  \vfill
  \noindent\textcolor{RuleGray}{\rule{\linewidth}{1pt}}\vspace{-2pt}
  \noindent\hfill\textit{Continued on next page}\par\vspace{-10pt}
  \noindent\textcolor{RuleGray}{\rule{\linewidth}{1pt}}
  \clearpage
}

% ---------- PDF hyperlinks ----------
\usepackage[hidelinks]{hyperref}
\hypersetup{
  colorlinks=false,
  pdfborder={0 0 0},
  linktoc=all,
  pdftitle={\DocTitle},
  pdfauthor={\ProgramName}
}

% =========================================================
\begin{document}

\section*{10.B — Adjustment Rules}
\phantomsection
\addcontentsline{toc}{section}{10.B — Adjustment Rules}

\subsection*{10.B.1 Measurement Inputs and Signal Quality}
\phantomsection
\addcontentsline{toc}{subsection}{10.B.1 Measurement Inputs and Signal Quality}

Minimum dataset aligned to Stage 2 logging posture (targets/rules only; schema remains Stage 2.E authority):

\begin{itemize}
\item \textbf{BODYCOMP}
\begin{itemize}
\item Weekly bodyweight average required (\textbf{ST-BW-WKAVG} / \textbf{MET-2B-001}); daily weigh-ins optional as inputs but not decision drivers.
\item Weekly circumference suite required (ST-CIRC-* / \textbf{MET-2B-003} through \textbf{MET-2B-009}): neck, chest, bicep, forearm, waist, thigh, calf.
\item Body fat estimate percent (\textbf{ST-BF-PCT} / \textbf{MET-2B-002}) is method-agnostic but trend-valid only when the same method/device and conditions are used across the comparison window.
\end{itemize}
\item \textbf{PERFORMANCE}/\textbf{QUALITY}
\begin{itemize}
\item Test-week outcomes when applicable (Stage 2 protocol outcomes under \textbf{VALID} test conditions).
\item Session RPE trend (\textbf{ST-RPE-0-10} / \textbf{MET-2B-042}) and qualitative performance notes (non-protocolized).
\end{itemize}
\item \textbf{RECOVERY}
\begin{itemize}
\item Sleep hours (\textbf{ST-SLEEP-HRS} / \textbf{MET-2B-038}) and sleep quality (\textbf{ST-SLEEP-QUAL-0-10} / \textbf{MET-2B-039}).
\item Soreness (\textbf{ST-SORENESS-0-10} / \textbf{MET-2B-040}) and fatigue (\textbf{ST-FATIGUE-0-10} / \textbf{MET-2B-041}).
\item Illness flags and systemic symptom notes (non-diagnostic).
\end{itemize}
\item \textbf{COMPLIANCE}
\begin{itemize}
\item On-plan day count (tracked operationally within nutrition logs; decision thresholds below).
\item Novelty events near gates: any meaningful diet structure change, new food class, new caffeine pattern, or new supplement category introduction inside protected windows.
\item Deload/Test-week flags (Stage 3 protected windows).
\end{itemize}
\end{itemize}

Signal quality rules (binding):

\begin{itemize}
\item Weekly bodyweight average drives decisions; single-day weigh-ins \textbf{MUST} NOT trigger structural changes.
\item Waist trend is evaluated only if tape method and landmarks are consistent; if inconsistent, treat waist trend as non-comparable and do not use it for plateau determination.
\item Minimum trend window before plateau decisions: at least 2 weeks of data unless safety triggers require earlier intervention.
\item No major changes from a single day or single weigh-in; trend-based posture only.
\item Do not interpret medical meaning from measurements; posture is operational and performance-support oriented.
\item Missing-data doctrine: if required fields are missing for the decision window, treat the signal as unreliable and default to \textbf{HOLD} posture until data quality is restored.
\end{itemize}

\newpage
\subsection*{10.B.2 Adjustment Cadence}
\phantomsection
\addcontentsline{toc}{subsection}{10.B.2 Adjustment Cadence}

Default cadence:

\begin{itemize}
\item Weekly review and decision point using the prior 7-day window for compliance and the last 14 days for plateau determination.
\item Hold periods after any change:
\begin{itemize}
\item 7–14 days hold unless safety requires earlier intervention (safety overrides timeline).
\end{itemize}
\item Test-week and gate-window restrictions:
\begin{itemize}
\item no novelty during Deload/Test weeks,
\item no major diet structure changes in the final 7–10 days before major gates,
\item during Deload/Test weeks, only distribution tweaks are allowed except the binding maintenance-calorie rule.
\end{itemize}
\end{itemize}

Required table:

\begingroup
\scriptsize
\setlength{\tabcolsep}{2pt}
\renewcommand{\arraystretch}{1.12}

\begin{longtblr}{
  width=\linewidth,
  rowhead=1,
  colspec = {
    X[l,1.75]
    X[l,2.70]
    X[l,2.40]
    X[l,2.85]
  },
  hlines,
  vlines,
  cells = {hsep=6pt, vsep=6pt, halign=l, valign=t},
  cell{1-Z}{1-4} = {fg=black},
  row{1} = {bg=StageBlue!15, fg=black, font=\bfseries, halign=l, valign=m},
  row{odd} = {bg=white},
  row{even} = {bg=LightGray},
}
\Hdr{Adjustment Window} &
\Hdr{Allowed Change Types} &
\Hdr{Not Allowed} &
\Hdr{Documentation Required} \\
Routine weekly review (non-protected weeks) &
Minor tweaks if $\geq$6/7 on-plan; structural changes if $\geq$12/14 on-plan; step-size calorie change $\pm$150–250 &
Single-day reactive changes; multiple simultaneous major changes &
Nutrition Adjustment Log entry: date, trigger evidence, change, hold period, review date, owner \\
7–10 days before major gate tests &
\textbf{HOLD} posture by default; only safety-required corrections; maintain stable foods &
Any major macro structure change; new food class; fiber overhaul; new supplement category &
Log any deviation as “novelty event” and treat as validity confounder; default to deferral \\
Deload/Test week (Stage 3 protected window) &
Calories set to maintenance; test-day carb redistribution of +50–100 g offset by reduced fat calories &
New restrictive cuts; new supplements; novelty foods; drastic fiber changes; caffeine pattern changes &
Log: maintenance-calorie posture applied; test-day carb offset executed; no novelty confirmed \\
Post-test week (immediate week after protected window) &
Resume prior posture if recovery stable; otherwise maintenance reset &
Aggressive deficit escalation in reaction to test results &
Log: decision, rationale, and hold period; treat recovery markers as primary constraint \\
\end{longtblr}

\endgroup

\newpage
\subsection*{10.B.3 Adjustment Levers and Priority Order}
\phantomsection
\addcontentsline{toc}{subsection}{10.B.3 Adjustment Levers and Priority Order}

Strict priority order (binding):

\begin{enumerate}
\item Compliance fixes before macro math
\begin{itemize}
\item Reduce friction first: simplify choices, reduce optional complexity, standardize meal timing and repeatable staples.
\item No meal plans generated; only operational simplification rules.
\end{itemize}
\item Calorie adjustment step (locked)
\begin{itemize}
\item Default structural calorie adjustment step: $\pm$150–250 kcal/day.
\item Only applied when compliance thresholds are met (see 10.B.4).
\end{itemize}
\item Macro distribution tweaks
\begin{itemize}
\item Protein floor remains stable within 10.A hybrid-anchor posture; do not cut protein to chase calories.
\item Fat floor remains 20–30\% except when offsetting test-day carb increase; do not undercut fat floor as a primary lever.
\item Carbs are the primary lever after protein and fat floor are satisfied.
\end{itemize}
\item Maintenance reset (locked)
\begin{itemize}
\item Maintenance reset equals 7 days at maintenance.
\item Used when recovery/performance collapses, diet fatigue risk rises, or restrictive behaviors appear.
\end{itemize}
\item Phase-specific bias
\begin{itemize}
\item Early phases tolerate deficit more; later phases default to maintenance/near-maintenance per 10.A; do not force loss targets when 10.A indicates maintenance/surplus posture.
\end{itemize}
\end{enumerate}

Required table:

\begingroup
\scriptsize
\setlength{\tabcolsep}{2pt}
\renewcommand{\arraystretch}{1.12}

\begin{longtblr}{
  width=\linewidth,
  rowhead=1,
  colspec = {
    X[l,1.55]
    X[l,2.55]
    X[l,2.15]
    X[l,2.75]
  },
  hlines,
  vlines,
  cells = {hsep=6pt, vsep=6pt, halign=l, valign=t},
  cell{1-Z}{1-4} = {fg=black},
  row{1} = {bg=StageBlue!15, fg=black, font=\bfseries, halign=l, valign=m},
  row{odd} = {bg=white},
  row{even} = {bg=LightGray},
}
\Hdr{Lever} &
\Hdr{Typical Use Case} &
\Hdr{Expected Effect} &
\Hdr{Risk / Control} \\
Compliance simplification &
Low on-plan days; planning failure; friction overload &
Improves adherence and restores signal quality &
Risk: overcomplication; control: simplify to staple posture and hold \\
Calorie step $\pm$150–250 &
Trend too slow (early phases) or too fast (any phase) under adequate compliance &
Adjusts weight trend without drastic disruption &
Risk: overcorrection; control: 7–14 day hold \\
Carb redistribution (no calorie change) &
Performance/recovery strain without bodyweight trend issue; test-week prep &
Supports training quality without increasing restriction &
Risk: GI novelty; control: familiar staples only near gates \\
Maintenance reset 7 days &
Sleep/performance collapse, diet fatigue, restrictive behavior risk &
Restores recovery markers and reduces validity contamination &
Risk: fear of “losing progress”; control: trend-based evaluation and defined exit rule \\
Phase bias enforcement &
Late-phase high consequence testing &
Prevents under-fueling that invalidates tests &
Risk: unwanted slow loss; control: accept stability to protect evidence \\
\end{longtblr}

\endgroup

\newpage
\subsection*{10.B.4 Bodyweight and Waist Trend Decision Rules}
\phantomsection
\addcontentsline{toc}{subsection}{10.B.4 Bodyweight and Waist Trend Decision Rules}

Locked bands and definitions apply. Compliance thresholds are gating conditions for structural changes.

Compliance thresholds (binding):

\begin{itemize}
\item Minor tweaks allowed at $\geq$6/7 on-plan days in the last 7 days.
\item Structural calorie changes require $\geq$12/14 on-plan days in the last 14 days.
\item If thresholds are not met: do not adjust calories; execute compliance simplification and \textbf{HOLD}.
\end{itemize}

Meaningful waist reduction (operational definition, non-medical):

\begin{itemize}
\item Meaningful waist reduction exists when either:
\begin{itemize}
\item the weekly waist trend decreases across two consecutive weeks using consistent tape method and consistent landmark posture, OR
\item the average of the last two weekly waist measures is lower than the prior two-week average by at least the measurement rounding unit used in Stage 2.E.
\end{itemize}
\item If method/landmark drift is present: treat waist signal as non-comparable and exclude from plateau logic.
\end{itemize}

Locked decision rules (if-then):

\begin{itemize}
\item Too fast loss (sustained):
\begin{itemize}
\item If $>0.75\%$ bodyweight/week sustained for 2 consecutive weeks AND compliance $\geq$12/14:
\begin{itemize}
\item raise calories by one step (+150 to +250 kcal/day),
\item preferentially raise carbs while preserving protein and fat floor,
\item hold 7–14 days unless safety dictates earlier adjustment.
\end{itemize}
\end{itemize}
\item Too slow loss (Phase 00–Phase 05 only):
\begin{itemize}
\item If $<0.10\%$ bodyweight/week sustained for 2 consecutive weeks AND compliance $\geq$12/14:
\begin{itemize}
\item reduce calories by one step (\verb|~|150 to \verb|~|250 kcal/day) OR tighten discretionary carbs while preserving protein and fat floor,
\item hold 7–14 days.
\end{itemize}
\item Phase 06+ note: later phases may be maintenance-focused; do not force loss targets when 10.A indicates maintenance/surplus posture.
\end{itemize}
\item Plateau (trend-based):
\begin{itemize}
\item If 14 days with $<0.10\%$ bodyweight change AND no meaningful waist reduction AND compliance $\geq$12/14:
\begin{itemize}
\item apply one calorie step reduction (\verb|~|150 to \verb|~|250 kcal/day) OR execute a structured compliance simplification pass first if compliance quality is borderline even though thresholds are met,
\item hold 7–14 days.
\end{itemize}
\end{itemize}
\item Bodyweight stable but waist improving:
\begin{itemize}
\item \textbf{HOLD}; do not cut further.
\end{itemize}
\item Bodyweight down but waist not improving and recovery poor:
\begin{itemize}
\item \textbf{HOLD} or execute maintenance reset; do not cut further.
\end{itemize}
\item Bodyweight up unexpectedly:
\begin{itemize}
\item assess sodium/carb swings, adherence drift, and novelty events,
\item \textbf{HOLD} 7 days unless the 2-week trend confirms upward drift under high compliance.
\end{itemize}
\end{itemize}

\begin{landscape}

Required decision table:

\begingroup
\scriptsize
\setlength{\tabcolsep}{2pt}
\renewcommand{\arraystretch}{1.12}

\begin{longtblr}{
  width=\linewidth,
  rowhead=1,
  colspec = {
    X[l,1.5]
    X[l,2.85]
    X[l,1.75]
    X[l,2.50]
    Q[l,wd=.65in]
    X[l,2.40]
  },
  hlines,
  vlines,
  cells = {hsep=2pt, vsep=6pt, halign=l, valign=t},
  row{odd} = {bg=white},
  row{even} = {bg=LightGray},
  row{1} = {bg=StageBlue!15, fg=black, font=\bfseries, halign=l, valign=m},
}
\Hdr{Scenario} &
\Hdr{Signal Definition trend-based} &
\Hdr{Immediate Action} &
\Hdr{Next-Step Adjustment} &
\Hdr{Hold Period} &
\Hdr{Notes} \\
Too fast loss &
$>0.75\%$ \textbf{BW}/week for 2 consecutive weeks; compliance $\geq$12/14 &
Increase calories +150–250 &
Prefer carbs upshift; maintain protein; fat within floor band &
7–14 days &
Protect recovery and validity; do not “reward” with novelty foods \\
Too slow loss (Phase 00–05) &
$<0.10\%$ \textbf{BW}/week for 2 consecutive weeks; compliance $\geq$12/14 &
Reduce calories by 150-250 &
Prefer carb reduction; preserve protein; keep fat floor &
7–14 days &
If compliance is borderline, simplify first; do not cut on poor compliance \\
Plateau &
14 days $<0.10\%$ \textbf{BW} change + no meaningful waist reduction; compliance $\geq$12/14 &
Simplify or step-reduce &
One step reduction if simplification already optimized &
7–14 days &
Waist signal used only if method/landmarks stable \\
\textbf{BW} stable, waist improving &
Waist down across 2 weeks under consistent method; \textbf{BW} flat &
\textbf{HOLD} &
None &
7–14 days &
Do not cut further; preserve performance/recovery \\
\textbf{BW} down, waist flat, recovery poor &
\textbf{BW} down but waist not improving; sleep/fatigue/soreness trending worse &
\textbf{HOLD} or maintenance reset &
7 days maintenance reset &
7 days reset, then 7-day reassess &
Avoid compounding restriction with recovery collapse \\
\textbf{BW} up, short-term &
\textbf{BW} up within 7 days; no 2-week upward trend; possible sodium/carb swing &
\textbf{HOLD} &
None unless 2-week trend confirms &
7 days &
Do not react to single week noise; check novelty events and compliance quality \\
Compliance below threshold &
$<12/14$ on-plan days &
Compliance simplification &
No structural calorie change allowed &
7 days &
Restore signal quality before adjustment \\
\end{longtblr}

\endgroup

\end{landscape}

\newpage
\subsection*{10.B.5 Performance and Recovery Protection Rules}
\phantomsection
\addcontentsline{toc}{subsection}{10.B.5 Performance and Recovery Protection Rules}

Do-not-cut and/or shift-to-maintenance triggers (non-medical, operational) using Stage 2 readiness markers:

Triggers must include:

\begin{itemize}
\item persistent sleep debt trend (sleep hours down and/or sleep quality down),
\item persistent elevated soreness or fatigue trend,
\item performance regression on key markers or failed test weeks under \textbf{VALID} evidence posture,
\item frequent \textbf{MODIFIED} or \textbf{INVALID} sessions due to fatigue or illness posture (Stage 1.F validity discipline),
\item illness flags and systemic symptoms.
\end{itemize}

Required table:

\begingroup
\scriptsize
\setlength{\tabcolsep}{2pt}
\renewcommand{\arraystretch}{1.12}

\begin{longtblr}{
  width=\linewidth,
  rowhead=1,
  colspec = {
    X[l,1.35]
    X[l,2.55]
    X[l,2.35]
    X[l,2.35]
    X[l,1.95]
  },
  hlines,
  vlines,
  cells = {hsep=2pt, vsep=6pt, halign=l, valign=t},
  row{odd} = {bg=white},
  row{even} = {bg=LightGray},
  row{1} = {bg=StageBlue!15, fg=black, font=\bfseries, halign=l, valign=m},
}
\Hdr{Trigger Cluster} &
\Hdr{Example Signals} &
\Hdr{Nutrition Response} &
\Hdr{Gate Impact} &
\Hdr{Documentation} \\
Sleep debt cluster &
\textbf{ST-SLEEP-HRS} (\textbf{MET-2B-038}) down for several days; \textbf{ST-SLEEP-QUAL-0-10} (\textbf{MET-2B-039}) low trend &
\textbf{HOLD} deficit; move to maintenance if persistent; consider maintenance reset &
Near gates: \textbf{HOLD} deficit and protect validity; reslot if tests would be compromised &
Log signals; log response; note hold/reset period \\
Fatigue/soreness cluster &
\textbf{ST-FATIGUE-0-10} (\textbf{MET-2B-041}) elevated trend; \textbf{ST-SORENESS-0-10} (\textbf{MET-2B-040}) elevated trend &
\textbf{HOLD} or maintenance reset; do not cut further &
Gate posture: avoid new restriction near protected windows; prioritize admissible evidence &
Log cluster; record decision; set review date \\
Valid performance regression &
\textbf{VALID} test regression vs required tier posture OR repeated poor session performance notes under high compliance &
\textbf{HOLD} deficit; consider maintenance reset; keep protein stable &
Gate posture: default to \textbf{HOLD} for progression if performance evidence degrades; protect retest feasibility &
Log test IDs and results; log nutrition action and hold \\
Validity contamination &
Increased \textbf{MODIFIED}/\textbf{INVALID} sessions due to fatigue/illness; inability to meet Cardio continuity reliably &
\textbf{HOLD} deficit; simplify; consider maintenance reset &
Gate posture: inadmissible evidence cannot support advancement; reslot per Stage 3 &
Log validity distribution and reasons; log nutrition response \\
Illness/systemic symptoms &
fever, acute respiratory illness, severe GI distress (non-diagnostic) &
Suspend deficit; move to maintenance; focus on hydration safety; resume only when stable &
Gate posture: tests reslot; \textbf{MEDICAL} \textbf{HOLD} may apply per governance when indicated &
Log symptom flags; log nutrition posture change; document return criteria \\
\end{longtblr}

\endgroup

Rules (binding):

\begin{itemize}
\item Near gates, recovery protection overrides weight-loss preference: default posture is hold maintenance and protect validity rather than cutting further.
\item No medical diagnosis; refer to escalation posture when red flags occur.
\end{itemize}

\newpage
\subsection*{10.B.6 Deload/Test Week Nutrition Adjustment Rules}
\phantomsection
\addcontentsline{toc}{subsection}{10.B.6 Deload/Test Week Nutrition Adjustment Rules}

Binding rule:

\begin{itemize}
\item On Deload/Test weeks, calories are set to maintenance (even if currently surplus posture exists in later phases, maintain maintenance posture for that Deload/Test week unless Stage 10.A explicitly already at maintenance; in either case, do not cut).
\end{itemize}

Allowed changes:

\begin{itemize}
\item On test day(s): +50–100 g carbs offset by equal calories reduced from fats.
\item Protein maintained within 10.A posture.
\item Fat remains within 20–30\% floor band after offset (the offset is a distribution change, not a floor violation).
\end{itemize}

Not allowed:

\begin{itemize}
\item new restrictive cuts,
\item new supplements or new supplement category introductions (7.03–7.06 changes discouraged; 7.02 never required),
\item novelty foods,
\item drastic fiber changes.
\end{itemize}

Required table:

\begingroup
\scriptsize
\setlength{\tabcolsep}{2pt}
\renewcommand{\arraystretch}{1.12}

\begin{longtblr}{
  width=\linewidth,
  rowhead=1,
  colspec = {
    X[l,1.55]
    X[l,2.10]
    X[l,2.75]
    X[l,2.45]
    X[l,2.10]
  },
  hlines,
  vlines,
  cells = {hsep=2pt, vsep=6pt, halign=l, valign=t},
  row{odd} = {bg=white},
  row{even} = {bg=LightGray},
  row{1} = {bg=StageBlue!15, fg=black, font=\bfseries, halign=l, valign=m},
}
\Hdr{Week Type} &
\Hdr{Goal} &
\Hdr{Allowed Changes} &
\Hdr{Not Allowed} &
\Hdr{Notes} \\
Deload/Test week &
Protect recovery and test validity &
Calories to maintenance; protein stable; carbs are upshifted; fat is reduced only to offset the test-day carb increase &
New deficits; new foods; new supplement category changes; caffeine pattern changes &
Treat as controlled-variable week; no novelty posture is binding \\
Test day within Deload/Test week &
Reduce fatigue contamination and stabilize output &
+50–100 g carbs; offset calories from fats; protein stable &
New GI-risk foods; high-fiber experiments; “performance hacks” &
Use familiar staples; log any deviation as novelty event \\
Non-test days in Deload/Test week &
Preserve stability &
Hold maintenance calories; stable foods &
“Catch-up restriction” after prior week &
Maintain stable routine; do not compensate \\
\end{longtblr}

\endgroup

\newpage
\subsection*{10.B.7 Life-Constraint Overrides}
\phantomsection
\addcontentsline{toc}{subsection}{10.B.7 Life-Constraint Overrides}

Minimum viable execution rules (operational, non-programming):

\begin{itemize}
\item Unstable access (food availability, schedule volatility)
\begin{itemize}
\item prioritize protein adequacy and calorie control using staples,
\item simplify food variety to reduce planning friction,
\item maintain tracking and on-plan day counting discipline.
\end{itemize}
\item Limited cooking tools
\begin{itemize}
\item operate within 7.13 minimal tooling posture,
\item assembly-meal posture allowed (simple combinations), no meal plan generation,
\item focus on repeatability and food safety.
\end{itemize}
\item Travel and work disruption
\begin{itemize}
\item portable baseline concept: stable “default meal structure” posture using accessible staples,
\item maintain tracking and on-plan day posture,
\item avoid novelty near gates; stabilize earlier.
\end{itemize}
\item Adherence collapse
\begin{itemize}
\item simplify to Tier 1 execution posture,
\item remove optional complexity and hold structural changes until compliance meets thresholds.
\end{itemize}
\end{itemize}

Required decision table:

\begingroup
\scriptsize
\setlength{\tabcolsep}{2pt}
\renewcommand{\arraystretch}{1.12}

\begin{longtblr}{
  width=\linewidth,
  rowhead=1,
  colspec = {
    X[l,1.20]
    X[l,1.85]
    X[l,2.20]
    X[l,2.85]
    X[l,1.55]
  },
  hlines,
  vlines,
  cells = {hsep=2pt, vsep=6pt, halign=l, valign=t},
  row{odd} = {bg=white},
  row{even} = {bg=LightGray},
  row{1} = {bg=StageBlue!15, fg=black, font=\bfseries, halign=l, valign=m},
}
\Hdr{Constraint} &
\Hdr{Example} &
\Hdr{Immediate Action} &
\Hdr{7-Day Plan operational} &
\Hdr{When to Re-evaluate} \\
Unstable access &
missed prep, irregular meal timing &
Simplify to staples; keep protein stable &
Repeatable staple pattern; track daily; maintain food safety &
Next weekly review \\
Limited tools &
no stove access temporarily &
Use 7.13 minimal tools; choose safe ready-to-eat staples &
Maintain protein and calories within posture; no novelty near gates &
7 days \\
Travel disruption &
hotel/work travel week &
Portable baseline; maintain tracking &
Hold macro posture; avoid new foods near gates; prioritize hydration safety &
Weekly review or post-travel \\
Adherence collapse &
$<6/7$ on-plan days &
Hold macro changes; simplify &
Remove optional foods/complexity; rebuild compliance to $\geq$12/14 for structural change &
After 7 days, then 14-day threshold \\
\end{longtblr}

\endgroup

\newpage
\subsection*{10.B.8 Documentation and Change Control}
\phantomsection
\addcontentsline{toc}{subsection}{10.B.8 Documentation and Change Control}

Nutrition Adjustment Log template:

\begingroup
\scriptsize
\setlength{\tabcolsep}{2pt}
\renewcommand{\arraystretch}{1.12}

\begin{longtblr}{
  width=\linewidth,
  rowhead=1,
  colspec = {
    X[l,1.55]
    X[l,2.35]
    X[l,2.10]
    X[l,4.00]
    X[l,3.40]
    X[l,2.35]
    X[l,2.10]
    X[l,2.10]
    X[l,6.00]
  },
  hlines,
  vlines,
  cells = {hsep=2pt, vsep=6pt, halign=l, valign=t},
  row{odd} = {bg=white},
  row{even} = {bg=LightGray},
  row{1} = {bg=StageBlue!15, fg=black, font=\bfseries, halign=l, valign=m},
}
\Hdr{Date} &
\Hdr{Phase/Week} &
\Hdr{Change} &
\Hdr{Trigger(s)} &
\Hdr{Intended Outcome} &
\Hdr{Hold Period} &
\Hdr{Review Date} &
\Hdr{Owner} &
\Hdr{Notes} \\
\end{longtblr}

\endgroup

Rules (binding):

\begin{itemize}
\item Owner must be one of: Coach Lead, Trainee, Medical, Equipment Lead.
\item Every change must include:
\begin{itemize}
\item trigger evidence summary (weekly \textbf{BW} avg + waist trend if comparable + recovery markers + compliance counts),
\item hold period (7–14 days unless safety),
\item review date.
\end{itemize}
\item No mid-gate novelty posture:
\begin{itemize}
\item changes inside the no-novelty window require explicit documentation and should default to deferral unless safety.
\end{itemize}
\item Changes must remain within 10.A posture; no independent diet system is permitted.
\end{itemize}

\newpage
\subsection*{10.B.9 Safety and Escalation Notes}
\phantomsection
\addcontentsline{toc}{subsection}{10.B.9 Safety and Escalation Notes}

\begin{itemize}
\item Red-flag symptoms override the plan:
\begin{itemize}
\item dizziness or fainting,
\item chest pain or pressure,
\item syncope or near-syncope,
\item severe GI distress,
\item heat illness signs,
\item \textrightarrow{} stop and seek evaluation per governance; resume only after stability is restored.
\end{itemize}
\item Disordered-eating warning signs (operational language):
\begin{itemize}
\item escalating restriction beyond rules,
\item binge/purge behaviors,
\item obsessional tracking distress that destabilizes execution,
\item avoidance of eating or persistent fear-driven under-eating,
\item \textrightarrow{} seek professional support posture; do not intensify restriction; default to maintenance reset and simplification.
\end{itemize}
\item Caffeine change posture (locked):
\begin{itemize}
\item no caffeine pattern changes within 14 days of gate windows,
\item outside that window, changes are allowed only with logging (date, reason, expected effect) and hold periods; do not stack caffeine changes with other major changes.
\end{itemize}
\end{itemize}

\newpage
\subsection*{10.B.10 Adjustment Integrity Checks}
\phantomsection
\addcontentsline{toc}{subsection}{10.B.10 Adjustment Integrity Checks}

\begin{itemize}
\item Decisions are trend-based, not single-day; weekly bodyweight average is the primary weight signal.
\item Calorie adjustment step size is enforced: $\pm$150–250 kcal/day; no larger jumps except via maintenance reset (planned 7-day maintenance).
\item Compliance thresholds are enforced:
\begin{itemize}
\item $\geq$6/7 on-plan days for minor tweaks,
\item $\geq$12/14 on-plan days for structural calorie changes.
\end{itemize}
\item Protein floor is maintained within 10.A hybrid-anchor posture; protein is not cut to “make room” for restriction.
\item Fat floor remains within 20–30\% of calories except for test-day carb offset mechanics; no chronic undercutting of fat floor.
\item Deload/Test weeks set calories to maintenance by rule; test-day carb redistribution (+50–100 g) is implemented by offsetting fat calories; no novelty foods.
\item No novelty near gates is enforced:
\begin{itemize}
\item no major structure changes in final 7–10 days before major gates,
\item caffeine pattern changes prohibited within 14 days of gate windows.
\end{itemize}
\item Every change is documented with date (YYYY-MM-DD), trigger evidence, hold period, review date, and owner role taxonomy.
\item Maintenance reset is used as the safety valve when recovery/performance collapses or restrictive behavior risk rises; exit only when readiness markers stabilize.
\item Feasibility under start-from-zero constraints is preserved; simplification precedes math; no meal plans or specialty procurement requirements.
\item Missing-data posture is enforced: \textbf{MISSING}/\textbf{UNKNOWN} recorded; no inference; default \textbf{HOLD} when signal quality is insufficient.
\item Any upstream conflict is labeled "Requires Stage 8 review" rather than silently changed; this deliverable does not rewrite 10.A.
\end{itemize}

\end{document}